% element: 10-s015:e04
\BeleskeID{10-s015}
Više oblika jednačine namerno je prikazano zato što pojedini rezultati mogu biti zbunjujući. Nastavak grafika započinje od granice prethodne oblasti, $V_I=V_{\gamma T}$. Kada u aktivni izraz ubacimo $0.5\,\mathrm{V}$, a zadržimo aktivni pad $0.6\,\mathrm{V}$, dobijamo $V_O>V_{CC}$.

Ovaj rezultat proizlazi iz diskretne prirode modela i naglog prelaska modelnog napona baza–emitor sa $0.5\,\mathrm{V}$ na $0.6\,\mathrm{V}$. Takve nelogičnosti prate grube modele; smatraćemo da su greške male i nastaviti sa njima. U crtanju se zanemaruje prekid na $V_I=V_{\gamma T}$, dok aktivni izraz zadržava svoj modelni $V_{BE}$. Kontinualna skica i numerički afini izraz zato služe različitim nivoima aproksimacije; ekstrapolacija jednačine van aktivne oblasti nije stvarno povećanje napona iznad napajanja.
